%------------------------------------------------------------------------------%
% Luis A. D'Afonseca
%
%------------------------------------------------------------------------------%
\documentclass[a4paper,12pt,fleqn]{article}

\usepackage{style_avaliacao}

\ShowAnswers

%------------------------------------------------------------------------------%
\begin{document}

\makeHeader{Integrais e Séries}{Prova Substitutiva -- Turma 5}{14/12/2023}

\vspace{\baselineskip}
\textbf{Substitui a prova} \underline{\hspace{20mm}}
\vspace{\baselineskip}

Os exercícios 1 e 2 valem 30 para quem perdeu a prova 1 ou 2 \\[-1mm]
O exercício 3 vale 40 para quem perdeu a prova 3 \\[-1mm]
O exercício 4 vale 40 para quem perdeu a prova 4

%------------------------------------------------------------------------------%
\vspace{\baselineskip}
\exercise{20}
Calcule a integral \(\int_0^1 \frac{\ln(x)}{\sqrt{x}} dx\).
\textit{Atenção ao intervalo de integração}.

\begin{answer}
  Calculando a primitiva
  \(
    F(x) = \int \frac{\ln(x)}{\sqrt{x}} dx
  \)
  por integral por partes
  \[
    u = \ln(x)
    \qquad
    dv = \frac{dx}{\sqrt{x}} = x^{-\sfrac{1}{2}}dx
  \]
  \[
    du = \frac{dx}{x}
    \qquad
    v = \frac{x^{\sfrac{1}{2}}}{\sfrac{1}{2}} = 2\sqrt{x}
  \]
  \begin{align*}
    F(x)
    & = \int \frac{\ln(x)}{\sqrt{x}} dx \\[2mm]
    & = \ln(x) 2 \sqrt{x} - \int 2\sqrt{x} \frac{dx}{x} \\[2mm]
    & = 2\ln(x)\sqrt{x} - 2\int x^{\sfrac{1}{2}-1} dx \\[2mm]
    & = 2\ln(x)\sqrt{x} - 2\int x^{-\sfrac{1}{2}} dx \\[2mm]
    & = 2\ln(x)\sqrt{x} - 2\sqrt{x} + C \\[2mm]
    & = 2\sqrt{x}\big( \ln(x) - 1\big) + C
  \end{align*}
  Para calcular a integral definida precisamos observar que temos uma assintota
  vertical em $x=0$ portanto essa é uma integral imprópria
  \begin{align*}
    I
    & = \int_0^1 \frac{\ln(x)}{\sqrt{x}} dx \\[2mm]
    & = \lim_{t\to0} \int_t^1 \frac{\ln(x)}{\sqrt{x}} dx \\[2mm]
    & = \lim_{t\to0} \left[ 2\sqrt{x}\big( \ln(x) - 1\big)\evalat{t}{1}\right]  \\[2mm]
    & = \lim_{t\to0} \left[ 2\sqrt{1}\big( \ln(1) - 1\big)
                          - 2\sqrt{t}\big( \ln(t) - 1\big) \right]   \\[2mm]
    & = 2( 0 - 1) -2 \lim_{t\to0} \left[\sqrt{t}\big(\ln(t) - 1\big) \right] \\[2mm]
    & = -2 -2 \lim_{t\to0} \left[\sqrt{t}\big(\ln(t) - 1\big) \right]
  \end{align*}
  Precisamos calcular o limite
  \begin{align*}
    L
    & = \lim_{t\to0} \sqrt{t}\big(\ln(t) - 1\big)  \\[2mm]
    & = \lim_{t\to0} \frac{\ln(t)-1}{t^{-\sfrac{1}{2}}} & \text{usando l'Hopital}\\[2mm]
    & = \lim_{t\to0} \frac{t^{-1}}{-\sfrac{1}{2}t^{-\sfrac{3}{2}}} \\[2mm]
    & = \lim_{t\to0} -2 t^{\sfrac{3}{2}-1} \\[2mm]
    & = -2 \lim_{t\to0} t^{\sfrac{1}{2}} \\[2mm]
    & = -2 \lim_{t\to0} \sqrt{t} = 0
  \end{align*}
  Portando
  \[
    I = \int_0^1 \frac{\ln(x)}{\sqrt{x}} dx = -2
  \]
\end{answer}

% 02
%------------------------------------------------------------------------------%
\vspace{\baselineskip}
\exercise{20}
Encontre o volume do sólido obtido pela rotação, em torno no eixo $y$,
da região contida entre as curvas \(y=x^3\), \(y=8\) e \(x=0\).

\begin{answer}
  Integrando por seções transversais em $y$ temos
  \[
    V = \int_a^b A(y) dy
  \]
  onde \(a=0\), \(b=8\) e
  \[
    A(y) = \pi r^2 = \pi \left(x^{\sfrac{1}{3}}\right)^2 = \pi x^{\sfrac{2}{3}}
  \]
  Assim
  \begin{align*}
      V
      & = \int_0^8 \pi x^{\sfrac{2}{3}} dy \\[2mm]
      & = \pi \int_0^8 x^{\sfrac{2}{3}} dy \\[2mm]
      & = \pi \frac{x^{\sfrac{5}{3}}}{\sfrac{5}{3}}\evalat{0}{8} \\[2mm]
      & = \frac{3\pi}{5} \left(8^{\sfrac{5}{3}}-0^{\sfrac{5}{3}}\right) \\[2mm]
      & = \frac{3\pi 2^5}{5} \\[2mm]
      & = \frac{96\pi}{5}
  \end{align*}
\end{answer}

% 03
%------------------------------------------------------------------------------%
\vspace{\baselineskip}
\exercise{20}
Analise a convergência da série \(\;\sum_{n=1}^\infty \frac{1}{2+3^n}\)

\begin{answer}
  Vamos comparar a série com a série geométrica convergente
  \[
    \sum_{n=1}^\infty \left(\frac{1}{3}\right)^n
  \]
  Primeiro verificamos que
  \[
    \frac{1}{2+3^n}
    < \frac{1}{3^n}
    = \left(\frac{1}{3}\right)^n
  \]
  Portanto, a série converge pelo teste da comparação.
\end{answer}

% 04
%------------------------------------------------------------------------------%
\vspace{\baselineskip}
\exercise{20}
Sabendo que \(f^{(n)}(x) = \frac{(-1)^nn!}{x^{n+1}}\), para \(n=0, 1, 2, \dots\)
\begin{enumerate}[label=\alph*)]
  \item Construa o Polinômio de Taylor de grau $n$ centrado em 2 da função $f$
  \item Determine o intervalo de convergência da Série de Taylor
  \item Determine o grau necessário para que a diferença entre o polinômio e a função
  seja menor do que $10^{-6}$ no intervalo $(2, 3)$
\end{enumerate}

\begin{answer}
\begin{enumerate}[label=\alph*)]
  \item O Polinômio de grau $n$ é
  \[
    P_n(x) = \sum_{k=0}^n \frac{f^{(k)}(2)}{k!}\left(x-2\right)^k
  \]
  como
  \[
    f^{(k)}(2)
    = \frac{(-1)^kk!}{x^{k+1}}
    = \frac{(-1)^kk!}{2^{k+1}}
  \]
  temos
  \begin{align*}
    P_n(x)
    & = \sum_{k=0}^n \frac{f^{(k)}(2)}{k!}\left(x-2\right)^k \\[2mm]
    & = \sum_{k=0}^n \frac{(-1)^kk!}{2^{k+1}}\frac{1}{k!}\left(x-2\right)^k \\[2mm]
    & = \sum_{k=0}^n \frac{(-1)^k}{2^{k+1}}\left(x-2\right)^k
  \end{align*}

  \item Podemos escrever a Série de Taylor omo uma série geométrica
  \[
    T(x)
    = \sum_{k=0}^\infty \frac{(-1)^k}{2^{k+1}}\left(x-2\right)^k
    = \frac{1}{2}\sum_{k=0}^\infty \left(\frac{2-x}{2}\right)^k
  \]
  que converge quando
  \begin{align*}
    \abs{r} = \abs{\frac{2-x}{2}} & < 1 \\[2mm]
    \abs{2-x}                     & < 2 \\[2mm]
    -2 < 2-x                      & < 2 \\[2mm]
    -4 < -x                       & < 0 \\[2mm]
    0  < x                        & < 4
  \end{align*}
  Portanto o intervalo de convergência é \((0, 4)\)

  \item O erro ao aproximar a função pelo Polinômio de Taylor de grau $n$
  é dada pelo Teorema de Taylor
  \[
    R_n(x)
    = \frac{f^{(n+1)}(\xi)}{(n+1)!}(x-2)^{n+1}
  \]
  para algum $\xi$ entre 2 e $x$. Neste caso particular
  \[
    R_n(x)
    = \frac{(-1)^{n+1}(n+1)!}{\xi^{n+2}}\frac{1}{(n+1)!}(x-2)^{n+1}
    = (-1)^{n+1}\frac{(x-2)^{n+1}}{\xi^{n+2}}
  \]
  portanto,
  \[
    \abs{R_n(x)}
    = \frac{\abs{x-2}^{n+1}}{\abs{\xi}^{n+2}}
  \]
  Como $2 \leq \xi \leq x < 3$, temos
  \[
    \frac{1}{\abs{\xi}^{n+2}}
    = \frac{1}{\xi^{n+2}}
    \leq \frac{1}{2^{n+2}}
  \]
  além disso,
  \[
    \abs{x-2}^{n+1}
    = (x-2)^{n+1}
    \leq 1^{n+1} = 1
  \]
  Assim
  \[
    \abs{R_n(x)}
    \leq \frac{1}{2^{n+2}}
  \]
  Para garantir que o erro seja menor do que a tolerância impomos
  \begin{align*}
    \frac{1}{2^{n+2}} & < 10^{-6}   & n + 2 & > 6\,\frac{\ln(10)}{\ln(2)}                  \\
    2^{n+2}           & > 10^6      & n     & > 6\,\frac{\ln(10)}{\ln(2)} - 2 \approx 17,9 \\
    (n+2)\ln(2)       & > 6\ln(10)
  \end{align*}

\end{enumerate}
\end{answer}

%------------------------------------------------------------------------------%
\end{document}
%------------------------------------------------------------------------------%
